\section{Case Study Setup}

The study runs on a single synthetic Trois-Rivi\`eres cold day with a minimum
ambient temperature of $-21.9\,^{\circ}\mathrm{C}$, discretized at 15-minute
resolution into $T = 96$ steps. We coordinate $N = 30$ electric-heating homes
with comfort band $\alpha = 0.5$, comfort weight $\lambda = 0.15$, flattening
weight $\mu = 1.0$, and ADMM penalty $\rho_{\mathrm{ADMM}} = 1.0$. The
communication-failure fraction is swept over
$\varrho \in \{0.0, 0.2, 0.4, 0.6, 0.8\}$.

Energy cost is evaluated against a time-of-use tariff with an off-peak rate of
0.078, a mid-peak rate of 0.128, and an evening-peak rate of 0.158~CAD/kWh.
Physical feasibility uses the IEEE 33-bus feeder with all loads at a power factor
of 0.95, with the aggregate scaled so the uncoordinated winter peak reaches the
feeder and trunk ampacity calibrated to a 115\% uncoordinated overload. The
software stack is pandapower~3.4.0 for power flow, scikit-learn~1.6.1 for the
cross-validation harness, and LightGBM~4.6.0 for the imputer. Every input,
intermediate, and reported figure is emitted as a governed artifact so the run
is reproducible from its declarative configuration.
Table~\ref{tab:params} collects the principal parameters.

\begin{table}[t]
\centering
\caption{Principal case-study parameters.}
\label{tab:params}
\begin{tabular}{@{}ll@{}}
\toprule
Parameter & Value \\
\midrule
Homes $N$ & 30 \\
Horizon / resolution & 96 steps / 15 min \\
Coldest-day min.\ temperature & $-21.9\,^{\circ}\mathrm{C}$ \\
Comfort band $\alpha$ & 0.5 \\
Comfort weight $\lambda$ & 0.15 \\
Flattening weight $\mu$ & 1.0 \\
ADMM penalty $\rho_{\mathrm{ADMM}}$ & 1.0 \\
Non-responsive fraction $\varrho$ & $\{0.0, 0.2, 0.4, 0.6, 0.8\}$ \\
TOU off / mid / peak (CAD/kWh) & 0.078 / 0.128 / 0.158 \\
Feeder & IEEE 33-bus, 12.66 kV \\
Power factor & 0.95 \\
Uncoordinated worst-line loading & 115\% (calibrated) \\
Power flow & AC Newton-Raphson \\
Software & pandapower 3.4.0, scikit-learn 1.6.1, LightGBM 4.6.0 \\
\bottomrule
\end{tabular}
\end{table}

\subsection{Reproducibility and governed artifacts}
The study is fully declarative. A \texttt{project.yaml} (a \emph{StudyProject}
contract) and a \texttt{workflow.yaml} (an eight-stage directed acyclic graph)
drive the entire pipeline; no step depends on hand-edited intermediate state.
The runner executes the stages in topological order --- synthetic cold-day
weather and per-home baseline generation, ADMM coordination, LightGBM imputer
training and cross-validation, the communication-failure sweep, power-flow
validation, and reporting --- and each stage emits a governed JSON report
following a fixed schema. Every report carries the same envelope fields
(\texttt{report\_id}, \texttt{schema\_version}, \texttt{created\_at},
\texttt{source\_domain}, \texttt{inputs}, \texttt{artifacts}, \texttt{summary},
\texttt{validation}), and each recorded artifact stores its byte count and SHA-256
digest, so every number in this paper is traceable to a specific file and run.
A regression baseline pins the headline aggregate peaks, and the run is
deterministic: fixed random seeds and a synthetic weather day remove any
external data dependency, so a reader can re-run the study from its configuration
and obtain baseline-matching results. We claim re-runnable, citable
reproducibility on this single configuration --- not statistical generality
across feeders or climates, which the limitations discuss. Table~\ref{tab:stages}
lists the eight stages and their principal governed outputs.

\begin{table}[t]
\centering
\caption{Workflow stages and principal governed outputs.}
\label{tab:stages}
\begin{tabular}{@{}rll@{}}
\toprule
\# & Stage & Principal output \\
\midrule
1 & Weather generation      & Synthetic cold-day temperature \\
2 & Baseline heating        & Per-home $h_{i,t}$ profiles \\
3 & Background load         & Fixed non-deferrable aggregate $B$ \\
4 & ADMM coordination       & Ideal $x_{i,t}$, iteration count \\
5 & Imputer training        & LightGBM model + CV metrics \\
6 & Communication sweep     & Per-$\varrho$ imputed aggregates \\
7 & Power-flow validation   & $V_{\min}$, line loading, verdict \\
8 & Reporting               & Figures, KPI tables, JSON report \\
\bottomrule
\end{tabular}
\end{table}
